\section{Training and serving Gleipnir monitors}
\label{sec:method}

\subsection{Monitor input and output}
\label{sec:interface}

A Gleipnir monitor reads one trajectory and returns one score.
The input is a single user message: an instruction that describes harmful behavior, followed by the trajectory between \texttt{<agent\_trajectory>} tags.
Appendix~\ref{app:prompts} gives the instructions in full, and Appendix~\ref{app:data-discrepancies} describes what the trajectories contain.
We apply the chat template of the model with thinking disabled and start the reply of the model with the text \texttt{Prediction:}.
The monitor then reads the logits $z_0$ and $z_1$ of the tokens \texttt{0} (benign) and \texttt{1} (harmful) at the next position.
The score is the probability of \texttt{1} among these two tokens,
\begin{equation}
s = \frac{e^{z_1}}{e^{z_0} + e^{z_1}} = \sigma(z_1 - z_0) .
\label{eq:score}
\end{equation}
This output has two properties that the output of Section~\ref{sec:background-sinha} lacks.
It needs one forward pass over the prompt and no generated text, so the monitor writes no rationale.
It is also continuous, so it ranks trajectories without the ties of an integer scale.

\paragraph{Effect of the output format.}
The output format changes the score of a monitor before any training.
We compare four formats on the unadapted Qwen3.5-9B, with the model, the OOD suite and the monitoring specification fixed; only the wording of the output instruction differs.
Binary logits score 0.698 Mean-OOD.
An integer from 0 to 10 scores 0.486.
Rounding the binary scores to one decimal lowers Mean-OOD by only 0.014, so the coarser scale does not explain this difference.
When the model writes its reasoning before the output, the integer scores 0.666 and the binary logits score 0.692, in both cases at a higher cost (Table~\ref{tab:ood-all}).
Binary logits without reasoning are therefore the cheapest of the four formats and score at least as high as the others.

\subsection{Soft-target distillation}
\label{sec:binary-logit-distillation}

We train the monitor to match the probability that a stronger teacher model gives to the decision token.
The teacher is Kimi K3.
It reads each training trajectory with the full specification (Appendix~\ref{app:prompts}) and answers with one decision token, without reasoning.
We read the log-probabilities of \texttt{0} and \texttt{1} for this token and normalize them as in Equation~\ref{eq:score}, which gives the teacher probability $p$.
We compute $p$ once for each of the 8,688 training trajectories and store it.
The loss of the monitor on a trajectory is the binary cross-entropy between its score $s$ and the teacher probability,
\begin{equation}
\mathcal{L} = -\,p \log s - (1 - p) \log (1 - s) .
\label{eq:loss}
\end{equation}
The loss has no term for the source label and no temperature.
In contrast to the pipeline of Section~\ref{sec:background-sinha}, the teacher writes no rationale, and no judge model or RL stage is needed.

\paragraph{Teacher-agreement filter.}
We use the source labels only to filter the training set.
We remove a trajectory when the decision of the teacher ($p \geq 0.5$) disagrees with its source label, which removes 852 of the 8,688 trajectories and leaves 7,836.
In a comparison on one seed of the 4B, the filter raised pAUROC@20 on the validation set from 0.847 to 0.896.
It also raised the false-positive rate at a threshold of 0.5 from 5.3\% to 21.4\%, so the filter improves the ranking of the monitor but not its calibration at this threshold.
This comparison has no control that removes the same number of trajectories at random.

\paragraph{Soft targets against hard labels.}
We test the effect of the soft targets by training on the source labels instead, with all other settings unchanged.
For the 4B, we compare three seeds for each objective.
With hard labels the seeds score 0.690 to 0.746 Mean-OOD, and with soft targets they score 0.777 to 0.813; the means are 0.716 and 0.792, a difference of 0.076.
On the validation set the means are 0.833 and 0.903.
For the 27B, we have one seed for each objective, and the difference is smaller.
With the compact instruction and BF16 serving, hard labels score 0.904 Mean-OOD and soft targets score 0.909.
In the configuration that we release (Sections~\ref{sec:training-setup} and~\ref{sec:serving}), they score 0.933 and 0.951.
On the validation set, the hard-label 27B scores higher in that configuration, 0.959 against 0.944.
Its scores are also less graded: 80\% of its OOD scores lie within 0.01 of 0 or 1, against 20\% for the soft-target 27B.
Soft targets therefore help the 4B on both sets, and for the 27B they help the OOD score and the spread of the scores but not the validation score.
One explanation is that soft targets add less when the unadapted model is already close to the teacher.
We cannot separate this from the size of the model with two models, and we leave it open.

\subsection{Training setup}
\label{sec:training-setup}

\todo{Data, LoRA configuration, the two instructions, seed averaging for the 4B, the instruction choice for the 27B, quantization-aware distillation.}

\subsection{Serving and cost}
\label{sec:serving}

\todo{Main serving decisions: single prefill pass, low-precision weights with QAD for the 4B, last-token final-layer attention, host-path changes. Cost per 1,000 evaluations at on-demand and committed-use B200 rates. Score change against BF16 serving. Released artifacts. Throughput per change goes in the appendix.}
